% preamble.tex -- 项目报告的共享导言区（由 Lead 维护，章节作者不要修改）
%
% 文档主类在 report.tex 中固定为 article + ctex：
%   \documentclass[11pt,a4paper]{article}
%   \usepackage[UTF8,fontset=fandol]{ctex}
% 因此本文件只放导言区设置，不放 \documentclass。

% ---------- 字体 ----------
% 按文件名加载 TeX Gyre（TeX Live 自带），不依赖系统字体名。
\setmainfont{texgyretermes}[
  Extension=.otf,
  UprightFont=*-regular,
  BoldFont=*-bold,
  ItalicFont=*-italic,
  BoldItalicFont=*-bolditalic]
\setsansfont{texgyreheros}[
  Extension=.otf,
  UprightFont=*-regular,
  BoldFont=*-bold,
  ItalicFont=*-italic,
  BoldItalicFont=*-bolditalic]
\setmonofont{texgyrecursor}[
  Extension=.otf,
  UprightFont=*-regular,
  BoldFont=*-bold,
  ItalicFont=*-italic,
  BoldItalicFont=*-bolditalic,
  Scale=0.9]

% ---------- 版面 ----------
\usepackage{geometry}
\geometry{top=2.4cm,bottom=2.4cm,left=2.6cm,right=2.6cm,headheight=15pt,footskip=1.0cm}
\setlength{\parindent}{2em}
\setlength{\parskip}{0.25em}
\setlength{\emergencystretch}{3em}
\renewcommand{\arraystretch}{1.15}

% 章节间距收紧（报告表格/图多，缩短标题前后留白）
\usepackage{titlesec}
\titlespacing*{\section}{0pt}{2.0ex plus 0.4ex minus 0.2ex}{1.0ex plus 0.2ex}
\titlespacing*{\subsection}{0pt}{1.5ex plus 0.3ex minus 0.2ex}{0.6ex plus 0.2ex}
\titlespacing*{\subsubsection}{0pt}{1.2ex plus 0.3ex minus 0.2ex}{0.5ex plus 0.2ex}

% 浮动体与正文的间距（表格密集，统一收紧）
\setlength{\textfloatsep}{10pt plus 2pt minus 2pt}
\setlength{\floatsep}{9pt plus 2pt minus 2pt}
\setlength{\intextsep}{9pt plus 2pt minus 2pt}
\setlength{\abovecaptionskip}{5pt}
\setlength{\belowcaptionskip}{2pt}

% ---------- 数学与表格 ----------
\usepackage{amsmath,amssymb}
\usepackage{booktabs}
\usepackage{tabularx}
\usepackage{longtable}
\usepackage{array}
\usepackage{multirow}
\usepackage{makecell}
\usepackage{enumitem}
\setlist{itemsep=0.2em,topsep=0.3em}

% ---------- 图形 ----------
\usepackage{graphicx}
\graphicspath{{figures/}}
\usepackage{tikz}
\usepackage{pgfplots}
\pgfplotsset{compat=1.18}
\usetikzlibrary{arrows.meta,positioning,shapes.geometric,fit,calc,backgrounds}
\usepackage{float}
\usepackage{caption}
\usepackage{subcaption}
\captionsetup{font=small,labelfont=bf,skip=5pt}

% ---------- 颜色与代码 ----------
\usepackage{xcolor}
\definecolor{mseblue}{RGB}{26,74,122}
\definecolor{msegray}{RGB}{95,99,104}
\definecolor{codebg}{gray}{0.96}
\definecolor{codeframe}{gray}{0.45}

\usepackage{listings}
\lstdefinestyle{cli}{%
  basicstyle=\ttfamily\footnotesize,
  breaklines=true,
  breakatwhitespace=true,
  columns=fullflexible,
  keepspaces=true,
  frame=single,
  framerule=0.3pt,
  rulecolor=\color{codeframe},
  backgroundcolor=\color{codebg},
  showstringspaces=false,
  tabsize=2,
  aboveskip=0.7em,
  belowskip=0.7em,
  xleftmargin=0.4em,
  xrightmargin=0.2em,
}
\lstdefinestyle{code}{%
  style=cli,
  language=C++,
  keywordstyle=\color{mseblue}\bfseries,
  commentstyle=\color{msegray}\itshape,
  stringstyle=\color{teal},
}

% 行内代码：\code{./build/mse build --stem}。
% \detokenize 保留空格与下划线等特殊字符；参数里不要出现字面百分号。
\newcommand{\code}[1]{\texttt{\detokenize{#1}}}

% ---------- 页眉页脚 ----------
\usepackage{fancyhdr}
\pagestyle{fancy}
\fancyhf{}
\fancyhead[L]{\small\color{msegray}迷你搜索引擎项目报告}
\fancyhead[R]{\small\color{msegray}2026 年 10 月}
\fancyfoot[C]{\small\thepage}
\renewcommand{\headrulewidth}{0.4pt}
\fancypagestyle{plain}[fancy]{%
  \fancyhf{}%
  \fancyfoot[C]{\small\thepage}%
  \renewcommand{\headrulewidth}{0pt}%
}

% ---------- 标题名（中文，显式设置以免依赖 ctex 默认值） ----------
\renewcommand{\figurename}{图}
\renewcommand{\tablename}{表}
\renewcommand{\contentsname}{目录}
\renewcommand{\listfigurename}{插图目录}
\renewcommand{\listtablename}{表格目录}
\renewcommand{\refname}{参考文献}
\renewcommand{\appendixname}{附录}

% ---------- 超链接 ----------
\usepackage[hidelinks,bookmarksnumbered=true]{hyperref}
\usepackage{url}

% ---------- 报告内约定宏 ----------
% 结论高亮：\finding{一句话结论}
\newcommand{\finding}[1]{%
  \par\medskip\noindent
  \begin{minipage}{\linewidth}%
    \colorbox{codebg}{\parbox{\dimexpr\linewidth-2\fboxsep}{%
      \noindent\textbf{\color{mseblue}结论：}#1}}%
  \end{minipage}%
  \par\medskip
}

% 未完成标记（终稿应为 0 处；用 \TODO{...} 可被 grep 出来）
\newcommand{\TODO}[1]{\textcolor{red}{\textbf{[TODO: #1]}}}

% 表格数字右对齐的简写列类型
\newcolumntype{N}{>{\raggedleft\arraybackslash}X}
\newcolumntype{R}{>{\raggedleft\arraybackslash}r}
